% Volume V: Admissibility and Continuation
% Part IX. A General Theory of Institutional State Change
% Chapter 47: Record, Verify, Recognize, Authorize
% Spine: proof-spines/ch047-spine.md

\chapter{Record, Verify, Recognize, Authorize}
\label{ch:ch047}

The foundational chain of \textbf{Record}, \textbf{Verify}, \textbf{Recognize}, \textbf{Authorize}, and \textbf{Continue} is:
\[
\operatorname{Record}
\centernot\Rightarrow
\operatorname{Verify}
\centernot\Rightarrow
\operatorname{Recognize}
\centernot\Rightarrow
\operatorname{Authorize}
\centernot\Rightarrow
\operatorname{Continue}.
\]
\Definition\ where each operator names a distinct institutional stage rather than an entailment.
Each term has its own formal role. Record creates persistence. Verification establishes specified relations between claim and trace. Recognition places the verified object within an institutional category. Authorization licenses an action. Continuation asks whether that authorization remains active across time.

The point of the chain is that no stage entails the next. A record may exist without verification; a verified trace may fail recognition; a recognized status may still not authorize coercion, payment, access, or speech; and an authorization may lapse, expire, or be interrupted rather than continue. Volume~V therefore begins from typed non-implication: institutional state change is trustworthy only when these operations remain distinct enough to be separately inspected.
